\section{Configuració de Sepolia i contractes desplegats}
\label{annex:sepolia}

Aquest annex recull la configuració blockchain utilitzada durant el desenvolupament i l’avaluació, així com les adreces dels contractes intel·ligents desplegats a la xarxa de proves Sepolia.

\subsection{Configuració de la xarxa}

La configuració general utilitzada és la següent:

\begin{table}[htbp]
\centering
\renewcommand{\arraystretch}{1.2}
\begin{tabularx}{\textwidth}{p{0.32\textwidth} X}
\hline
\textbf{Paràmetre} & \textbf{Valor} \\
\hline
Xarxa & Sepolia \\
Identificador de cadena & \texttt{11155111} \\
Identificador hexadecimal & \texttt{0xaa36a7} \\
Moneda nativa & Sepolia ETH \\
Proveïdor RPC & Google Cloud Web3 Portal \\
Explorador de blocs & Etherscan Sepolia \\
Entorn de desplegament & Hardhat \\
\hline
\end{tabularx}
\caption{Configuració de la xarxa utilitzada durant les proves}
\label{tab:configuracio_sepolia_annex}
\end{table}

L’URL concreta del node RPC no es publica a la memòria quan incorpora credencials o identificadors privats. Aquesta informació es proporciona mitjançant variables d’entorn durant l’execució.

\subsection{Contractes desplegats}

La Taula~\ref{tab:contractes_desplegats} recull els contractes utilitzats durant les proves.

\begin{table}[htbp]
\centering
\small
\renewcommand{\arraystretch}{1.2}
\begin{tabularx}{\textwidth}{
    p{0.28\textwidth}
    p{0.38\textwidth}
    X
}
\hline
\textbf{Contracte} &
\textbf{Adreça a Sepolia} &
\textbf{Finalitat} \\
\hline

\texttt{SecurityTestContract} &
\url{0xA32D1b3BE52922ffA5e65403d9378165c672a363} &
Generació d’aprovacions ERC20 limitades, revocacions i aprovacions elevades. \\

\texttt{SecurityTestNFT} &
\url{0x216033047ffbec7eb223bf206e4d440f644500d3} &
Generació de permisos globals ERC721 mitjançant \texttt{setApprovalForAll}. \\

\hline
\end{tabularx}
\caption{Contractes intel·ligents desplegats a Sepolia}
\label{tab:contractes_desplegats}
\end{table}

\subsection{Verificació dels contractes}

Els contractes es van verificar mitjançant Etherscan, fet que permet consultar públicament el codi font, l’ABI i les funcions exposades. Aquesta verificació també facilita la generació manual d’operacions mitjançant la interfície \textit{Write Contract}.

Per a cada contracte es pot registrar:

\begin{itemize}
    \item l’adreça de desplegament;
    \item el bloc de desplegament;
    \item el \textit{hash} de la transacció de desplegament;
    \item la versió del compilador Solidity;
    \item els paràmetres del constructor;
    \item l’estat de verificació a Etherscan.
\end{itemize}

\begin{table}[htbp]
\centering
\small
\renewcommand{\arraystretch}{1.2}
\begin{tabularx}{\textwidth}{p{0.27\textwidth} X X}
\hline
\textbf{Element} &
\textbf{SecurityTestContract} &
\textbf{SecurityTestNFT} \\
\hline
Adreça & \texttt{0xA32D...a363} & \texttt{0x2160...00d3} \\
Bloc de desplegament & No conservat & \texttt{10588546} \\
Hash del desplegament & No conservat & \texttt{0xbdd099cf...27417ccb} \\
Pragma de Solidity & \texttt{\string^0.8.20} & \texttt{\string^0.8.20} \\
Verificació & Verificat a Etherscan & Verificat a Etherscan \\
\hline
\end{tabularx}
\caption{Informació tècnica dels contractes desplegats}
\label{tab:informacio_contractes_annex}
\end{table}

La identitat del contracte \texttt{SecurityTestNFT} es va recuperar a partir de l'historial ERC721 del compte de proves i es va contrastar amb Etherscan. El contracte va ser creat per \texttt{0x8527...81aa}, el mateix compte emprat durant l'avaluació, i el rebut del desplegament presenta estat satisfactori. Aquestes dades permeten reproduir els casos D1 i D2 directament des de la interfície \textit{Write Contract} del contracte verificat.
